% ==============================================================================
% ПАРТИЯ B03-030: DISS-005-B1616 .. DISS-005-B1642
% Глава 3: Файл данных COMTRADE и унификация параметров
% Замечания: DISS-005-C0118 .. DISS-005-C0120 (3 шт.)
% Объекты: 0 шт.
% ==============================================================================

% DISS-005-B1616
\subsubsection{Файл данных (.dat)}\label{sec:ch3_dat_file}

% DISS-005-B1617
Возможны два варианта

\begin{itemize}[itemsep=2pt,topsep=2pt]
  % DISS-005-B1618
  \item ASCII: Текстовый формат, легко читаемый, но занимает больше места.
  % DISS-005-B1619
  \item BINARY: Двоичный формат (16-битные или 32-битные целые числа), более компактный и предпочтительный для хранения больших объемов данных.
\end{itemize}

% DISS-005-B1620
Исходный датасет содержит файлы в обоих форматах, так как первичная обработка не включала конвертацию типа .dat файла, а лишь замену его имени на MD5-хеш. Для дальнейшего использования и удобства использования, как писалось ранее, все файлы будут приведены к формату ASCII.\par\vspace{0.2cm}

% DISS-005-B1621
\subsection{Необходимость деперсонализации данных}\label{sec:ch3_depersonalization_necessity}

% DISS-005-B1622
Для публичного использования осциллограмм и их предоставления третьим лицам необходима деперсонализация (анонимизация) данных [92, 93, 94]. Без этого важного шага необработанные данные, используемые в исследованиях и часто содержащие служебную или конфиденциальную информацию, не могут быть переданы или процитированы. Это ограничение не только препятствует независимому доступу к исходным наборам данных, но и делает невозможной проверку моделей и результатов, тем самым подрывая прозрачность, воспроизводимость и достоверность исследований. Таким образом, анонимизация служит важным этапом для сбалансирования потребности в доступности данных с защитой конфиденциальной информации.\par\vspace{0.2cm}

% DISS-005-B1623
Отметим, что для внутренних исследований и анализа на предприятии исходная информация об объекте, времени и конкретном устройстве может быть чрезвычайно важна. Поэтому, наряду с созданием анонимизированного публичного датасета, сохраняется возможность работы с исходными, не деперсонализированными данными, используя ранее описанные функции поиска по конкретным терминалам (раздел 3.3.4). Однако для целей данной работы и формирования общедоступного ресурса все осциллограммы проходят процедуру деперсонализации.\par\vspace{0.2cm}

% DISS-005-B1624
\subsection{Процесс деперсонализации}\label{sec:ch3_depersonalization_process}

% DISS-005-B1625
Процесс деперсонализации реализован в виде набора функций (в основном, в классе ProcessingOscillograms, например, метод deleting\_\allowbreak confidential\_\allowbreak information\_\allowbreak in\_\allowbreak all\_\allowbreak files) и включает следующие основные операции, преимущественно с конфигурационными файлами .cfg формата COMTRADE:\par\vspace{0.2cm}

\begin{itemize}[itemsep=2pt,topsep=2pt]
  % DISS-005-B1626
  \item \textbf{Определение и унификация кодировки:} Файлы .cfg могут быть сохранены в различных кодировках (UTF-8, Windows-1251, OEM 866 и др.). На первом этапе осуществляется попытка корректного определения исходной кодировки, для чего была выработана эвристика последовательной проверки наиболее распространенных кодировок (utf-8 --\textgreater{} windows-1251 --\textgreater{} OEM 866), с последующим пересохранением файла в едином формате UTF-8 для обеспечения кроссплатформенной совместимости и корректного отображения символов.
  % DISS-005-B1627
  \item \textbf{Удаление идентификационной информации:}
  \begin{itemize}[itemsep=2pt,topsep=2pt]
    % DISS-005-B1628
    \item Из первой строки .cfg файла, содержащей, как правило, название объекта, номер регистратора и версию формата, удаляется информация об объекте и регистраторе.
    % DISS-005-B1629
    \item Даты и время начала и окончания записи осциллограммы, а также дата и время срабатывания триггера, заменяются на стандартизированное значение (например, 01/01/0001, 01:01:01.000000). Это делается как для основных временных меток события, так и для даты последней модификации файла.
  \end{itemize}
  % DISS-005-B1630
  \item \textbf{Анонимизация имен файлов:} Исходные имена файлов .cfg и соответствующих им .dat файлов заменяются на их MD5-хеш-сумму (рассчитанную по содержимому .dat файла). Это обеспечивает полную анонимность имен и сохраняет связь между конфигурационным файлом и файлом данных, при этом гарантируя уникальность имени для уникального содержимого.
  % DISS-005-B1631
  \item \textbf{Логирование проблемных файлов:} В процессе деперсонализации файлы, которые не удалось корректно прочитать из-за проблем с кодировкой, защитой от записи или поврежденной структурой, добавляются в специальный список (protect\allowbreak ed\_\allowbreak files.txt) для последующего анализа или ручной обработки.
\end{itemize}\par\vspace{0.2cm}

% DISS-005-B1632
\subsection{Унификация параметров осциллограмм}\label{sec:ch3_parameters_unification}

% DISS-005-B1633
После деперсонализации ключевой задачей становится унификация различных параметров осциллограмм для обеспечения их сопоставимости и возможности совместной обработки.\par\vspace{0.2cm}

% DISS-005-B1634
\subsubsection{Группировка по частоте сети и частоте дискретизации}\label{sec:ch3_frequency_grouping}

% DISS-005-B1635
Осциллограммы, записанные различными устройствами, могут соответствовать сетям с разной номинальной частотой (как правило, 50 Гц или 60 Гц) и иметь различные частоты дискретизации (например, 1200, 1600, 6400 выборок в секунду). Эти параметры существенно влияют на обработку данных, так как число выборок на период номинальной частоты (например, 20 мс для 50 Гц) напрямую зависит от частоты дискретизации.\par\vspace{0.2cm}

% DISS-005-B1636
Для систематизации данных была реализована функция («grouping\_\allowbreak by\_\allowbreak sampling\_\allowbreak rate\_\allowbreak and\_\allowbreak network» в «ProcessingOscillograms»), которая:\par\vspace{0.2cm}

\begin{itemize}[itemsep=2pt,topsep=2pt]
  % DISS-005-B1637
  \item Извлекает информацию о номинальной частоте сети и частоте дискретизации из «.cfg» файла (с помощью вспомогательной функции «extract\_frequencies»). Эта информация обычно содержится в строках после описания аналоговых и дискретных каналов.
  % DISS-005-B1638
  \item Перемещает файлы осциллограмм («.cfg» и «.dat») в поддиректории, названные в соответствии с их частотными характеристиками (например, «f\_\allowbreak network=50\_\allowbreak f\_\allowbreak rate=1600»).
  % DISS-005-B1639
  \item В случаях, когда не удается однозначно определить эти параметры (например, из-за нестандартного формата «.cfg» файла или отсутствия данных), файлам присваиваются условные значения (1/1), и они помещаются в специальную категорию для последующего анализа и ручной коррекции.
  % DISS-005-B1640
  \item Результаты такой группировки представлены в Таблице 2\commentnote{DISS-005-C0118}{Алексей Евдаков [2]}{Сделать автоматизированной ссылкой, пока что это не так}, где показано количество осциллограмм для различных комбинаций частот. Также в таблице приведена сопутствующая информация по обработке этих групп: по количеству не «пустых осциллограмм» (Раздел 3.6.4) и количеству осциллограмм у которых после всех унификаций оказались одинаковые имена сигналов (дубликаты) внутри одной осциллограммы (Раздел 3.4.4.2). Более подробные детали которой будут описаны далее.
\end{itemize}\par\vspace{0.2cm}

% DISS-005-B1641
\par\vspace{0.2cm}

% DISS-005-B1642
\paragraph{Таблица 2. Количество осциллограмм с различными частотными параметрами\commentnote{DISS-005-C0119}{Алексей Евдаков [2]}{Находится пока что тут, но может поменяться\newline\newline D:\textbackslash{}DataSet\textbackslash{}depersonalized\_ALL\_OSC\_v2.7\newline Файл\newline Статистика v1}\commentnote{DISS-005-C0120}{Алексей Евдаков [2]}{Можно и с экселя вставить, посмотрим как будет лучше.}}\label{tab:ch03_oscillograms_frequency_parameters}
\mbox{}\par\vspace{0.2cm}

% ==============================================================================
% ПАРТИЯ B03-031: DISS-005-B1643 .. DISS-005-B1799
% Глава 3: Таблица 2 (Частотные параметры осциллограмм, TAB-077, 157 ячеек)
% ==============================================================================

\mbox{}\par\vspace{0.2cm}
{\small
\setlength{\tabcolsep}{4.5pt}
\noindent
\begin{tabular}{|>{\centering\arraybackslash}p{32pt}|>{\centering\arraybackslash}p{36pt}|>{\centering\arraybackslash}p{86pt}|>{\centering\arraybackslash}p{68pt}|>{\centering\arraybackslash}p{51.5pt}|>{\centering\arraybackslash}p{68pt}|>{\centering\arraybackslash}p{51.5pt}|}
\hline
% DISS-005-B1643
\multicolumn{2}{|c|}{\textbf{Частота, Гц}} &
% DISS-005-B1644
\multirow{2}{*}{\textbf{\begin{tabular}[c]{@{}c@{}}Количество\\ осциллограмм\end{tabular}}} &
% DISS-005-B1645
\multicolumn{2}{c|}{\textbf{\begin{tabular}[c]{@{}c@{}}Не пустых\\ осциллограмм\end{tabular}}} &
% DISS-005-B1646
\multicolumn{2}{c|}{\textbf{\begin{tabular}[c]{@{}c@{}}Осциллограммы\\ с дубликатами\end{tabular}}} \\
\cline{1-2} \cline{4-7}
% DISS-005-B1647
\textbf{сети} &
% DISS-005-B1648
\textbf{АЦП} &
% DISS-005-B1649
&
% DISS-005-B1650
\textbf{Количество} &
% DISS-005-B1651
\textbf{\begin{tabular}[c]{@{}c@{}}Их\\ процент\end{tabular}} &
% DISS-005-B1652
\textbf{Количество} &
% DISS-005-B1653
\textbf{\begin{tabular}[c]{@{}c@{}}Их\\ процент\end{tabular}} \\
\hline
% DISS-005-B1654
\multirow{17}{*}{50} &
% DISS-005-B1655
1 &
% DISS-005-B1656
26 &
% DISS-005-B1657
- &
% DISS-005-B1658
x &
% DISS-005-B1659
3 &
% DISS-005-B1660
11,54\% \\ \cline{2-7}
% DISS-005-B1661
 &
% DISS-005-B1662
600 &
% DISS-005-B1663
666 &
% DISS-005-B1664
347 &
% DISS-005-B1665
52,10\% &
% DISS-005-B1666
45 &
% DISS-005-B1667
6,76\% \\ \cline{2-7}
% DISS-005-B1668
 &
% DISS-005-B1669
800 &
% DISS-005-B1670
23 &
% DISS-005-B1671
13 &
% DISS-005-B1672
56,52\% &
% DISS-005-B1673
16 &
% DISS-005-B1674
69,57\% \\ \cline{2-7}
% DISS-005-B1675
 &
% DISS-005-B1676
916 &
% DISS-005-B1677
1 &
% DISS-005-B1678
1 &
% DISS-005-B1679
100,00\% &
% DISS-005-B1680
 &
% DISS-005-B1681
 \\ \cline{2-7}
% DISS-005-B1682
 &
% DISS-005-B1683
1000 &
% DISS-005-B1684
925 &
% DISS-005-B1685
486 &
% DISS-005-B1686
52,54\% &
% DISS-005-B1687
217 &
% DISS-005-B1688
23,46\% \\ \cline{2-7}
% DISS-005-B1689
 &
% DISS-005-B1690
1100 &
% DISS-005-B1691
2~753 &
% DISS-005-B1692
1~841 &
% DISS-005-B1693
66,87\% &
% DISS-005-B1694
207 &
% DISS-005-B1695
7,52\% \\ \cline{2-7}
% DISS-005-B1696
 &
% DISS-005-B1697
1200 &
% DISS-005-B1698
30~666 &
% DISS-005-B1699
11~736 &
% DISS-005-B1700
38,27\% &
% DISS-005-B1701
9~777 &
% DISS-005-B1702
31,88\% \\ \cline{2-7}
% DISS-005-B1703
 &
% DISS-005-B1704
1600 &
% DISS-005-B1705
15~453 &
% DISS-005-B1706
6~290 &
% DISS-005-B1707
40,70\% &
% DISS-005-B1708
76 &
% DISS-005-B1709
0,49\% \\ \cline{2-7}
% DISS-005-B1710
 &
% DISS-005-B1711
1800 &
% DISS-005-B1712
90 &
% DISS-005-B1713
71 &
% DISS-005-B1714
78,89\% &
% DISS-005-B1715
11 &
% DISS-005-B1716
12,22\% \\ \cline{2-7}
% DISS-005-B1717
 &
% DISS-005-B1718
2000 &
% DISS-005-B1719
15 &
% DISS-005-B1720
8 &
% DISS-005-B1721
53,33\% &
% DISS-005-B1722
 &
% DISS-005-B1723
 \\ \cline{2-7}
% DISS-005-B1724
 &
% DISS-005-B1725
2100 &
% DISS-005-B1726
2 &
% DISS-005-B1727
2 &
% DISS-005-B1728
100,00\% &
% DISS-005-B1729
 &
% DISS-005-B1730
 \\ \cline{2-7}
% DISS-005-B1731
 &
% DISS-005-B1732
2400 &
% DISS-005-B1733
40 &
% DISS-005-B1734
21 &
% DISS-005-B1735
52,50\% &
% DISS-005-B1736
9 &
% DISS-005-B1737
22,50\% \\ \cline{2-7}
% DISS-005-B1738
 &
% DISS-005-B1739
4000 &
% DISS-005-B1740
2 &
% DISS-005-B1741
2 &
% DISS-005-B1742
100,00\% &
% DISS-005-B1743
 &
% DISS-005-B1744
 \\ \cline{2-7}
% DISS-005-B1745
 &
% DISS-005-B1746
4800 &
% DISS-005-B1747
45 &
% DISS-005-B1748
9 &
% DISS-005-B1749
20,00\% &
% DISS-005-B1750
 &
% DISS-005-B1751
 \\ \cline{2-7}
% DISS-005-B1752
 &
% DISS-005-B1753
6400 &
% DISS-005-B1754
15 &
% DISS-005-B1755
6 &
% DISS-005-B1756
40,00\% &
% DISS-005-B1757
7 &
% DISS-005-B1758
46,67\% \\ \cline{2-7}
% DISS-005-B1759
 &
% DISS-005-B1760
8000 &
% DISS-005-B1761
13 &
% DISS-005-B1762
10 &
% DISS-005-B1763
76,92\% &
% DISS-005-B1764
4 &
% DISS-005-B1765
30,77\% \\ \cline{2-7}
% DISS-005-B1766
 &
% DISS-005-B1767
12800 &
% DISS-005-B1768
2 &
% DISS-005-B1769
2 &
% DISS-005-B1770
100,00\% &
% DISS-005-B1771
 &
% DISS-005-B1772
 \\ \hline
% DISS-005-B1773
\multirow{3}{*}{60} &
% DISS-005-B1774
1200 &
% DISS-005-B1775
18 &
% DISS-005-B1776
18 &
% DISS-005-B1777
100,00\% &
% DISS-005-B1778
3 &
% DISS-005-B1779
16,67\% \\ \cline{2-7}
% DISS-005-B1780
 &
% DISS-005-B1781
1920 &
% DISS-005-B1782
6 &
% DISS-005-B1783
6 &
% DISS-005-B1784
100,00\% &
% DISS-005-B1785
 &
% DISS-005-B1786
 \\ \cline{2-7}
% DISS-005-B1787
 &
% DISS-005-B1788
3840 &
% DISS-005-B1789
4 &
% DISS-005-B1790
4 &
% DISS-005-B1791
100,00\% &
% DISS-005-B1792
 &
% DISS-005-B1793
 \\ \hline
% DISS-005-B1794
\multicolumn{2}{|c|}{Сумма} &
% DISS-005-B1795
50765 &
% DISS-005-B1796
20~873 &
% DISS-005-B1797
41,12\% &
% DISS-005-B1798
10~375 &
% DISS-005-B1799
20,44\% \\ \hline
\end{tabular}
}

